\chapter{Kawasaki Disease}
\index{Kawasaki Disease}

\section*{Definition and Overview}
Kawasaki disease is an acute illness that causes inflammation of the blood vessels (vasculitis) throughout the body, most commonly affecting children under five. It presents with fever lasting more than five days, a rash, red eyes, red cracked lips and a strawberry tongue, swollen lymph nodes, and swelling of the hands and feet. The most important complication is inflammation of the coronary arteries, which can cause aneurysms and, rarely, heart attacks in untreated children. Prompt treatment with immunoglobulin and aspirin markedly reduces this risk. The cause is unknown, but it is believed to be an abnormal immune response to an infection in genetically susceptible children.

\section*{Epidemiology}
Kawasaki disease is the leading cause of acquired heart disease in children in developed countries. It is most common in children under five, with a peak around 18 to 24 months, and is more common in boys and in children of Asian descent. In many countries, around 100 to 200 children are diagnosed each year, and the incidence is rising. Most cases occur in winter and spring, and although the cause is unknown, the pattern suggests an infectious trigger. With early treatment, the risk of serious heart complications is low.

\section*{Aetiology and Pathophysiology}
The cause of Kawasaki disease is unknown, but the evidence points to an abnormal immune response to one or more infections in genetically susceptible children, rather than a contagious disease itself. The immune system mounts an intense inflammatory response that damages the walls of small and medium blood vessels, particularly the coronary arteries that supply the heart. In the acute phase, children have high fever and systemic inflammation; in the subacute phase, the risk of coronary aneurysm is highest. Without treatment, up to one in four children develops coronary artery abnormalities. Treatment with intravenous immunoglobulin calms the inflammation and cuts the aneurysm risk to around 2 to 4 per cent.

\section*{Clinical Features}
\begin{itemize}
    \item Fever lasting five days or more, usually high and unresponsive to paracetamol
    \item A rash, often on the trunk and nappy area
    \item Red, bloodshot eyes without discharge
    \item Red, cracked lips, a strawberry tongue, and redness of the mouth and throat
    \item Swollen lymph nodes, usually on one side of the neck
    \item Swelling and redness of the hands and feet, with peeling of the fingers and toes in the second week
    \item Irritability, poor feeding, and tiredness
    \item Incomplete Kawasaki disease: some features absent, making diagnosis harder
\end{itemize}

\section*{Differential Diagnosis}
Many childhood illnesses resemble Kawasaki disease.
\begin{itemize}
    \item \textbf{Viral infections:} measles, adenovirus, and other viruses with fever and rash
    \item \textbf{Scarlet fever and streptococcal infection:} fever, rash, and strawberry tongue
    \item \textbf{Toxic shock syndrome and staphylococcal infection}
    \item \textbf{Juvenile idiopathic arthritis:} prolonged fever and rash
    \item \textbf{Drug reactions:} fever and rash from medicines
    \item \textbf{Other vasculitides:} rare inflammatory conditions
\end{itemize}

\section*{When to Seek Medical Care}
Any child with a fever lasting more than five days, especially with a rash, red eyes, cracked lips, or swollen hands and feet, should be assessed promptly, as early treatment of Kawasaki disease prevents heart complications. Parents should not wait for all features to appear. A child with high fever and any of the features should be seen by a doctor urgently.

\textbf{Contact your local emergency services for:}
\begin{itemize}
    \item A child who is very unwell, floppy, or difficult to wake
    \item Chest pain, breathlessness, or a racing heartbeat in a child with Kawasaki disease
    \item Seizures or convulsions
    \item Signs of dehydration: dry mouth, no wet nappies, or sunken eyes
    \item A rash that does not fade when pressed
\end{itemize}

\section*{Diagnosis and Investigations}
\begin{itemize}
    \item \textbf{Clinical criteria:} the diagnosis is based on the characteristic features and fever duration
    \item \textbf{Blood tests:} inflammatory markers, blood count, and liver function
    \item \textbf{Urine tests:} to exclude infection
    \item \textbf{Echocardiogram:} assessment of the coronary arteries, repeated over time
    \item \textbf{Electrocardiogram:} for heart rhythm and changes
    \item \textbf{Exclusion of other causes:} cultures and viral testing
    \item \textbf{Cardiology referral:} for children with or at risk of coronary changes
\end{itemize}

\section*{Management}
\begin{itemize}
    \item \textbf{Intravenous immunoglobulin:} given early, dramatically reducing coronary complications
    \item \textbf{Aspirin:} high-dose in the acute phase, then low-dose for several weeks
    \item \textbf{Hospital care:} monitoring of fever, hydration, and heart status
    \item \textbf{Retreatment:} for children who do not respond to the first immunoglobulin dose
    \item \textbf{Additional anti-inflammatory medicines:} for resistant disease
    \item \textbf{Cardiology follow-up:} echocardiograms for weeks to months
    \item \textbf{Long-term monitoring:} for children who develop coronary aneurysms
    \item \textbf{Family support:} education and reassurance
\end{itemize}

\section*{Complications}
\begin{itemize}
    \item Coronary artery inflammation and aneurysms
    \item Heart attack, which is rare but can occur
    \item Inflammation of the heart muscle or valves
    \item Long-term coronary narrowing requiring ongoing care
    \item Irritability, dehydration, and poor feeding in the acute illness
    \item Very rare, late complications in untreated disease
\end{itemize}

\section*{Prognosis and Outcomes}
With early treatment, the outlook for Kawasaki disease is excellent: more than 95 per cent of treated children recover fully, and the risk of serious coronary complications falls to around 2 to 4 per cent. Untreated, up to a quarter develop coronary abnormalities, so early recognition is critical. Most children return to normal health and activity within weeks. Those with coronary aneurysms need long-term cardiology care, and a small number require treatment for narrowed arteries.

\section*{Prevention and Self-Care}
\begin{itemize}
    \item Seek medical care promptly for a child with prolonged fever and the features of Kawasaki disease
    \item Do not wait for all the symptoms to appear before seeking help
    \item Follow the treatment and cardiology follow-up plans
    \item Give medicines as prescribed and attend all reviews
    \item Keep the child comfortable, hydrated, and rested
    \item Be aware that the cause is unknown and it is not contagious
    \item Attend the echocardiogram appointments, which are essential
    \item Seek urgent care for chest pain, breathlessness, or a very unwell child
\end{itemize}

\section*{Sources and Further Reading}
The following reputable sources provide further information for healthcare students and patients seeking reliable, current advice about Kawasaki disease.
\begin{itemize}
    \item Royal Children's Hospital Melbourne. \textit{Kawasaki disease information.} https://www.rch.org.au/
    \item Sydney Children's Hospitals Network. \textit{Kawasaki disease.} https://www.schn.health.nsw.gov.au/
    \item HeartKids Australia. \textit{Support for children with heart conditions.} https://heartkids.org.au/
    \item Healthdirect Australia. \textit{Kawasaki disease.} https://www.healthdirect.gov.au/
    \item NHS. \textit{Kawasaki disease.} https://www.nhs.uk/conditions/kawasaki-disease/
    \item American Heart Association. \textit{Kawasaki disease.} https://www.heart.org/
\end{itemize}
